\subsubsection{Energia cinética}
\label{sec:dinAcoplada_energia_cinetica}

A energia cinética total é escrita como a soma da contribuição associada à
atitude da espaçonave (tratada como corpo rígido) e das
contribuições de cada um dos $n$ elos do manipulador robótico, isto é,

\begin{equation}
T = T_{B} + \sum_{i=1}^{n} T_{i},
\label{eq:T_total}
\end{equation}

onde $n$ é a quantidade de elos do \gls{mr}, neste trabalho considera-se $n = 3$.

De forma geral, a energia cinética da espaçonave robótica (base + \gls{mr}) poderia ser decomposta em uma parcela translacional, associada ao movimento do seu centro de massa, e uma parcela rotacional, associada ao movimento de atitude em torno desse mesmo ponto.
Neste trabalho, a dinâmica de translação do centro de massa global é tratada através do movimento orbital relativo no referencial \gls{lvlh}, enquanto que o termo translacional associado ao centro de massa não é repetida na formulação lagrangiana da dinâmica acoplada. Portanto, define-se $T_B$ como a parcela rotacional da energia cinética da espaçonave em torno do seu centro de massa.

Quando somente a atitude é considerada como grau de liberdade dinâmico, essa
energia cinética de atitude é escrita como
\begin{equation}
T_{B} =
\frac{1}{2}\, {}^{B}\!\vec{\omega}_{B}^{\,T}\,\mathbf{I}_B(\vec{q})\,{}^{B}\!\vec{\omega}_{B},
\label{eq:TB}
\end{equation}
em que $\mathbf{I}_B(\vec{q})$ é o tensor de inércia da \emph{espaçonave como
um todo} (base + elos do manipulador) em torno do seu centro de massa global
$G$, expresso no referencial $\{B\}$ fixo ao veículo. Em particular,
$\mathbf{I}_B(\vec{q})$ depende da configuração articular $\vec{q}$, uma vez
que a redistribuição de massa provocada pelo movimento dos elos altera a
posição de $G$ e o tensor de inércia equivalente da espaçonave.

Na formulação acoplada adotada, os graus de liberdade de translação da base
não são incluídos no vetor de coordenadas generalizadas: assume-se que o
movimento do centro de massa global é prescrito pelo modelo de translação
orbital. Em um modelo completamente livre de forças externas, a conservação da
quantidade de movimento linear implicaria um pequeno deslizamento da
espaçonave ao longo da órbita em resposta ao deslocamento dos elos; tal efeito
é assumido de segunda ordem e permanece embutido no problema de translação, não
sendo retroalimentado explicitamente na dinâmica de atitude.

A energia cinética do $i$-ésimo elo do manipulador é composta pelas parcelas
translacional e rotacional usuais,
\begin{equation}
T_i =
\frac{1}{2} m_i\, {}^{B}\!\vec v_{G_i}^{\,T}\,{}^{B}\!\vec v_{G_i}
+
\frac{1}{2}\, {}^{B}\!\vec\omega_i^{\,T}\, I_i\, {}^{B}\!\vec\omega_i,
\label{eq:Ti}
\end{equation}
em que ${}^{B}\!\vec v_{G_i}$ e ${}^{B}\!\vec\omega_i$ são, respectivamente, a
velocidade linear do centro de massa e a velocidade angular do elo $i$,
expressas no referencial $\{B\}$ e obtidas a partir do Jacobiano acoplado
introduzido na Seção~\ref{sec:cinematica_diferencial_jacobiano}, e $I_i$ é o
tensor de inércia do elo $i$ em torno do seu próprio centro de massa.

Do ponto de vista físico, o termo translacional em \eqref{eq:Ti} representa o movimento do centro de massa do elo $i$ em relação ao corpo da espaçonave, por meio do quadrado da sua velocidade linear ${}^{B}\!\vec v_{G_i}$.
Em um sistema verdadeiramente livre de forças externas, a soma dessas contribuições está
associada a um deslocamento efetivo do centro de massa global da espaçonave, de modo que o movimento dos elos induz um pequeno deslizamento da nave na órbita para garantir a conservação da quantidade de movimento linear.
No presente trabalho, esse efeito é tratado no nível do problema de translação orbital e não é modelado explicitamente na dinâmica acoplada de atitude e juntas.


Agrupando-se as coordenadas generalizadas em um único vetor, define-se
\begin{equation}
\vec y
=
\begin{bmatrix}
\vec\eta_B \\
\vec q
\end{bmatrix},
\qquad
\dot{\vec y}
=
\begin{bmatrix}
\dot{\vec\eta}_B \\
\dot{\vec q}
\end{bmatrix}
\in \mathbb{R}^{3+n},
\label{eq:veloc_generalizadas}
\end{equation}
em que $\vec\eta_B \in \mathbb{R}^3$ reúne os parâmetros de atitude da base
(por exemplo, os ângulos de Euler ZYX da espaçonave) e $\vec q \in \mathbb{R}^n$
contém as coordenadas articulares do manipulador. Note-se que a velocidade
angular da base, ${}^{B}\!\vec\omega_B$, não é, em si, uma velocidade
generalizada, mas é obtida como combinação linear das velocidades
generalizadas de atitude, isto é,
${}^{B}\!\vec\omega_B = \mathbf{T}_\omega(\vec\eta_B)\,\dot{\vec\eta}_B$,
conforme a relação usual entre taxas de ângulos de Euler e velocidade angular.

Com essa definição, a energia cinética total pode ser escrita de forma compacta como
\begin{equation}
T =
\frac{1}{2}\,\dot{\vec y}^{\,T}\,
M(\vec y)\,
\dot{\vec y},
\label{eq:T_quadratica_Mq}
\end{equation}

em que $M(\vec y)\in\mathbb{R}^{(3+n)\times(3+n)}$ é a matriz de inércia
generalizada do conjunto base--manipulador, função da configuração
$(\vec\eta_B,\vec q)$, construída a partir das contribuições da base e dos
elos segundo as expressões \eqref{eq:TB} e \eqref{eq:Ti}.